\section{Introduction}
\label{sec:intro}

When a large language model produces a false statement, what went wrong? The answer matters enormously for safety: an LLM that hallucinates because it lacks knowledge is qualitatively different from one that ``knows'' the correct answer but reports something else under social pressure. Yet current lie-detection benchmarks and methods treat all false outputs as a single undifferentiated category~\citep{azaria2023internal, lin2021truthfulqa}, making it impossible to know which failure mode a detector is actually identifying---or whether it will generalize to the failure mode we care about most.

{\bf Why distinguishing falsehood mechanisms matters.} The growing deployment of LLMs in high-stakes domains---medical advice, legal reasoning, educational tutoring---demands that we understand not just \emph{whether} a model is wrong, but \emph{why}. A model that confabulates due to genuine uncertainty can be addressed with better calibration or retrieval augmentation. A model that systematically repeats misconceptions from its training data requires targeted debiasing. But a model that overrides its own knowledge to agree with a user---sycophancy---poses a fundamentally different and more concerning risk, one that existing detectors may miss entirely.

{\bf The gap in existing work.} Prior methods for detecting LLM falsehoods fall into three largely disconnected threads. Internal-state probing methods such as SAPLMA~\citep{azaria2023internal} and CCS~\citep{burns2022discovering} train classifiers on hidden activations but do not categorize \emph{why} a statement is false. Semantic entropy~\citep{farquhar2024detecting} specifically targets confabulations but by design cannot detect high-confidence incentive-driven errors. Sycophancy research~\citep{sharma2024towards} measures behavioral outcomes but does not quantify how sycophancy-driven errors compare in prevalence to epistemic failures. The HACK framework~\citep{simhi2025hack} introduces a Knowledge $\times$ Certainty taxonomy but does not examine incentive-driven misreporting. No existing study measures the prevalence breakdown of different falsehood types using a unified methodology.

{\bf Our approach.} We propose \ours, a behavioral classification framework that operationally distinguishes three mechanisms behind LLM false outputs: \emph{confabulations} (inconsistent wrong answers reflecting epistemic uncertainty), \emph{incentive-driven errors} (correct answers that flip under social pressure), and \emph{systematic errors} (consistently wrong answers reflecting memorized misconceptions). For each question, we collect multiple responses under baseline and social-pressure conditions, then classify errors by their consistency and pressure sensitivity. This approach requires no internal model access, making it applicable to any API-served LLM.

We evaluate \ours on 900 questions from \syceval~\citep{sharma2024towards} and \truthfulqa~\citep{lin2021truthfulqa} across two state-of-the-art models. Our key finding is that approximately 61\% of LLM errors are confabulations, 18--23\% are incentive-driven, and 15--21\% are systematic---and this distribution is remarkably stable across models ($\chi^2 = 1.23$, $p = 0.54$). A text-based classifier achieves 71.9\% accuracy at distinguishing these types, 38.6 percentage points above chance.

We make the following contributions:
\begin{itemize}[leftmargin=*,itemsep=0pt,topsep=0pt]
    \item We propose \ours, a behavioral framework that classifies LLM false outputs into three mechanism categories using multi-sample consistency and paired pressure testing, without requiring internal model access.
    \item We conduct the first prevalence measurement of falsehood mechanisms across two production LLMs, finding that ${\sim}20\%$ of errors are incentive-driven---a statistically significant fraction for \gemini ($p = 0.016$).
    \item We demonstrate that the three falsehood types produce partially distinguishable text signatures (71.9\% cross-validated accuracy), confirming they are distinct phenomena that warrant separate evaluation in lie-detection benchmarks.
\end{itemize}
